% =====================================================================
%  Section: Why Hopf Fibrations Belong in Particle Physics
%  Paper 1, "It from Bit via Gödel"
%  Drafted 2026-06-08.
%
%  This is a \input fragment for the main document. It assumes:
%    - biblatex is loaded and \autocite is available;
%    - the companion file hopf_particle_physics_refs_2026-06-08.bib
%      has been added to the bibliography resources;
%    - an amsthm-style `proposition' environment exists in the preamble.
%
%  Internal cross-references use placeholder labels:
%    sec:self-reference   -- the self-reference -> complex amplitudes chain
%    sec:division-algebras-- the Cayley--Dickson / Hurwitz step
%    sec:fanout-direction -- the FANOUT / measurement direction-selection argument
%    sec:entanglement-classes -- the FOLLOWING section (GHZ->quarks, W->leptons)
%  Re-point these to the real labels before compiling.
% =====================================================================

\section{Why Hopf Fibrations Belong in Particle Physics}
\label{sec:hopf-justification}

Invoking Hopf fibrations in a discussion of the Standard Model invites an immediate and fair objection: that the move is decorative, an analogy between two pieces of attractive mathematics with no load-bearing connection between them. The purpose of this section is to dispose of that objection before the construction is put to work. The connection is not an analogy. It is carried by two independent and independently published theorems, and the structures it relies upon are not selected for convenience but compelled by a pair of classification results. We take as already established the correspondence between qubit entanglement and the Hopf fibrations themselves --- the subject of a mature literature \autocite{mosseri2001hopf,bernevigchen2003geometry} --- and do not reproduce it here. What this section addresses is the prior question of why that correspondence is entitled to carry physical content at all.

\subsection{Two theorems, joined end to end}

The bridge from qubit entanglement to the Standard Model is built from two distinct results, neither of which is original to the present framework and neither of which depends on the other.

The first is geometric. The unit sphere in the Hilbert space of $n$ qubits is the total space of a Hopf fibration: $S^3 \to S^2$ for one qubit, $S^7 \to S^4$ for two, and $S^{15} \to S^8$ for three. For two and three qubits the fibration is \emph{entanglement-sensitive} --- separability is precisely the condition under which the base degenerates to a product of single-qubit Bloch spheres and the additional dimensions associated with the higher imaginary directions collapse \autocite{mosseri2001hopf,bernevigchen2003geometry}. This is a statement about the geometry of state spaces and requires no physical interpretation whatever.

The second is algebraic. Singling out a preferred complex structure on the octonions --- equivalently, fixing the split $\mathbb{O} \cong \mathbb{C} \oplus \mathbb{C}^3$ --- leaves a residual symmetry, and that residual symmetry is exactly the gauge group of the Standard Model. This was established through the exceptional Jordan algebra by Dubois-Violette and Todorov \autocite{duboisviolette2016exceptional,duboisviolettetodorov2019exceptional,todorovduboisviolette2018deducing} and given its sharpest spinorial form by Krasnov \autocite{krasnov2021so9}, who showed that the relevant group is the stabiliser of a complex structure inside $\mathrm{Spin}(9)$. We record the result in the form in which it enters the present work.

\begin{proposition}[Dubois-Violette--Todorov; Krasnov]\label{prop:gsm}
Let $J_R$ be the complex structure on the space $\mathbb{O}^2$ of $\mathrm{Spin}(9)$ spinors determined by a choice of unit imaginary octonion, equivalently the split $\mathbb{O} \cong \mathbb{C} \oplus \mathbb{C}^3$. The subgroup of $\mathrm{Spin}(9)$ commuting with $J_R$ is
\[
  G_{\mathrm{SM}} \;=\; \frac{SU(3) \times SU(2) \times U(1)}{\mathbb{Z}_6},
\]
acting on the singled-out $\mathbb{C}$ summand as a colour singlet of hypercharge $-1$ and on the $\mathbb{C}^3$ summand as a colour triplet of hypercharge $\tfrac{1}{3}$, both as $SU(2)$ doublets --- the quantum numbers of one left-handed generation.
\end{proposition}

The two results are independent, and the join between them is the only place a physical theory must do work. Each result is a theorem; what neither supplies is a reason why nature should fix one complex direction within $\mathbb{O}$ rather than treating all directions symmetrically. We return to this --- the genuine load-bearing assumption --- in \S\ref{sec:hopf-honest-debt}, and resolve it, in candidate form, in \S\ref{sec:fanout-direction}.

\subsection{The ladder is forced, not chosen}

It might still be objected that selecting the octonions, or the Hopf fibrations, is itself a modelling choice carrying as much latitude as the choice of a grand-unified group. It is not, and the reason is a pair of classification theorems that close off the alternatives entirely.

By Hurwitz's theorem the only normed division algebras are $\mathbb{R}$, $\mathbb{C}$, $\mathbb{H}$ and $\mathbb{O}$ \autocite{hurwitz1898composition,baez2002octonions}. By Adams's resolution of the Hopf-invariant-one problem --- together with the parallelisability results of Bott--Milnor and Kervaire --- the only parallelisable spheres are $S^0$, $S^1$, $S^3$ and $S^7$, so that there exist exactly three nontrivial Hopf fibrations, in bijection with $\mathbb{C}$, $\mathbb{H}$ and $\mathbb{O}$ \autocite{adams1960hopf,bottmilnor1958parallelizability,kervaire1958nonparallelizability}. These are not heuristics; they are among the more durable theorems in topology and algebra. Their force here is that, once a framework is committed to amplitudes valued in a normed division algebra --- a commitment that in the present construction is made far upstream, by the self-referential argument of \S\ref{sec:self-reference} and \S\ref{sec:division-algebras}, and not by fiat at this stage --- the ladder of Hopf fibrations is not selected but compelled, and it terminates of its own accord.

\begin{proposition}[Termination]\label{prop:termination}
The Cayley--Dickson tower of normed division algebras and the tower of parallelisable spheres both terminate at real dimension eight. Consequently there are exactly three nontrivial Hopf fibrations, the largest being the octonionic $S^{15} \to S^8$, and the construction admits no extension to a fourth case.
\end{proposition}

The termination is worth dwelling on, because it is unusual for a model of this kind to predict its own boundary. The next rung of the Cayley--Dickson sequence, the sedenions, possesses zero divisors; the Hopf map can no longer be globally defined; there is no fourth fibration. This is not a limitation imposed from outside but a structural fact, and it carries a physical consequence that we develop later: there is no room in the construction for a fourth gauge force or a fifth interaction, so that the discovery of either would count as evidence against the framework rather than as an invitation to extend it. The endpoint is a theorem, not a stipulation.

\subsection{Independent arrival as corroboration}

The most persuasive evidence that the destination is not an artefact of the route is that it is reached by routes sharing none of the same machinery. Szangolies \autocite{szangolies2025standardmodel} reaches $G_{\mathrm{SM}}$ from the emergent-gravity and Kaluza--Klein programmes, reading the singled-out complex direction as a dimensional reduction of a $9+1$-dimensional spacetime; the present framework reaches it from self-reference (\S\ref{sec:self-reference}); and Furey and Gillard--Gresnigt reach the same gauge group, and the same three-generation structure, through the minimal left ideals of the Dixon and Clifford algebras, without invoking Hopf fibrations at all \autocite{furey2018three,gillardgresnigt2019three}; the relation between their minimal ideals and the entanglement classes used here is taken up in \S\ref{sec:canonical-orbits}. The octonionic route to particle physics is in any case an old one, going back to the use of $\mathbb{O}$ to model the strong force \autocite{gunaydingursey1974quark}.

That programmes with disjoint starting points converge on $SU(3) \times SU(2) \times U(1)/\mathbb{Z}_6$ and on the same quark--lepton division is the coincidence Szangolies himself remarks upon: the characteristic structure of the Standard Model, routinely regarded as arbitrary, turns up independently in unrelated settings, and its appearing twice would then be doubly odd unless the two settings are in fact related. We read the convergence the same way --- as evidence that the division-algebraic description is singled out by the mathematics rather than imposed by the modeller, and therefore that Proposition~\ref{prop:gsm} is structural rather than contrived.

\subsection{The honest debt}
\label{sec:hopf-honest-debt}

It would be misleading to leave the impression that this route is free of assumptions, so we state plainly the one assumption that all of these programmes share. Proposition~\ref{prop:gsm} requires that \emph{some} complex direction in $\mathbb{O}$ be singled out, and the uniqueness of that choice is not supplied by the theorem. Krasnov parametrises the complex structure $J_R$ by an arbitrary unit imaginary octonion, and the residual group is the same for every such choice because $G_2 = \mathrm{Aut}(\mathbb{O})$ acts transitively on the unit imaginary octonions. The theorem therefore guarantees the gauge group \emph{given} that a direction is fixed, while remaining silent on \emph{why} one is fixed. Szangolies treats this as a free choice and flags it, in his own conclusions, as an open element of the construction --- one to which several inequivalent answers are available.

This, and not the appeal to Hopf fibrations, is the genuine load-bearing assumption of the whole approach. The Hopf fibrations are forced (Proposition~\ref{prop:termination}); the singling-out is not. The present framework offers a candidate resolution that the prior literature lacks: the act of measurement, modelled as \textsc{fanout}, selects the direction physically, and the transitivity of $G_2$ then guarantees that the gauge group is independent of which direction the measurement happens to select. This argument is given in \S\ref{sec:fanout-direction}; we flag here, at the appropriate epistemic status, that it is structural rather than fully proved, and that the residual question --- why a measurement selects \emph{exactly one} imaginary direction rather than merely \emph{permitting} a selection --- remains open.

A second debt is shared with every programme of this character and is not closed by anything said here: why the abstract quantum state should admit a tensor-product decomposition into qubits in the first place. This is the quantum factorisation, or mereology, problem. The entanglement-sensitivity on which the construction turns is meaningful only once such a decomposition is fixed, and we treat the decomposition as a necessary input rather than a derived result. Naming it is preferable to letting it pass unremarked.

\subsection{What this sets up}

The entanglement-sensitivity of the Hopf maps is not a side feature but the hinge of the physical reading that follows. Because separability is exactly the collapse of the base to a product of Bloch spheres, the structure of the entanglement --- the SLOCC class of a multi-qubit state, and within it the canonical maximally entangled orbit --- is carried by precisely the degrees of freedom on which $G_{\mathrm{SM}}$ acts in Proposition~\ref{prop:gsm}. This is what makes it possible, in \S\ref{sec:entanglement-classes}, to read the two genuinely tripartite entanglement classes onto the $\mathbb{C} \oplus \mathbb{C}^3$ split through their canonical representatives: the GHZ orbit confined to the colour-triplet $\mathbb{C}^3$ sector, the W orbit reaching the $\mathbb{C}$ singlet, recovering the quark--lepton division not as an assignment but as a consequence of the entanglement structure. The scope of those statements --- canonical orbits rather than whole classes, and the identification of the canonical orbits with the physical particles --- is made precise there. The Hopf fibrations are the bridge; the entanglement classes, treated next, are the traffic that crosses it.
